\documentclass[11pt]{article}

\newcommand{\cnum}{Math 245B}
\newcommand{\ced}{Winter 2026}
\newcommand{\ctitle}[4]{\title{\vspace{-0.5in}\cnum, \ced\\Assignment #1: #2}\author{\vspace{-0.35in}\\#3, #4}}
\usepackage{enumitem}
\usepackage[usenames,dvipsnames,svgnames,table,hyperref]{xcolor}
\usepackage{amsmath, amsfonts}
\usepackage{graphicx} % Required for inserting images
\usepackage{float}
\usepackage{listings}
\usepackage{xcolor}
\usepackage{hyperref}
\usepackage{comment}

\renewcommand*{\theenumi}{\alph{enumi}}
\renewcommand*\labelenumi{(\theenumi)}
\renewcommand*{\theenumii}{\roman{enumii}}
\renewcommand*\labelenumii{\theenumii.}

\definecolor{codegreen}{rgb}{0,0.6,0}
\definecolor{codegray}{rgb}{0.5,0.5,0.5}
\definecolor{codepurple}{rgb}{0.58,0,0.82}
\definecolor{backcolour}{rgb}{0.95,0.95,0.92}
\definecolor{header}{HTML}{205459}
\definecolor{solution}{HTML}{204d52}

\newcommand{\solution}[1]{{{\textcolor{header}{Solution:} \textcolor{solution}{#1}}}}
\setlist[enumerate]{label=(\alph*)}

\lstdefinestyle{mystyle}{
    backgroundcolor=\color{backcolour},
    commentstyle=\color{codegreen},
    keywordstyle=\color{magenta},
    numberstyle=\tiny\color{codegray},
    stringstyle=\color{codepurple},
    basicstyle=\ttfamily\footnotesize,
    breakatwhitespace=false,
    breaklines=true,
    captionpos=b,
    keepspaces=true,
    numbers=left,
    numbersep=5pt,
    showspaces=false,
    showstringspaces=false,
    showtabs=false,
    tabsize=2
}


\begin{document}
\ctitle{\#1}{}{Asher Christian}{006-150-286}
\date{}
\maketitle
\[
    \mathbb{E} \sum_{i_1,\dots,i_k}^{}X_{i_1}\dots X_{i_k} \text{ or } X \rightarrow Y
\] 
\[
    =\sum_{\pi \in P(k)}^{}\sum_{j_1,\dots,j_{|\pi|}}^{} \mathbb{E} X_{i_1}\dots X_{i_k}|_{\pi, J}
\] 
In this problem $X_i = 0/1$ and $Y_i = 0 / 1$ so
 \[
     X_{i_1}\dots X_{i_k} |_{\pi,J} = Y_{i_1}\dots Y_{i_k} |_{\pi, J}
\] 
We need
\[
\mathbb{E}(\sum_{i}^{}Y_i)^{k} \le C_k
\] 
\[
\sum_{\pi }^{}1 \le D_k \qquad D_k \in \mathbb{R}
\] 
if we fix $\pi$
\[
\sum_{J}^{}1 \le n^{|\pi|}
\] 
and
\[
    \mathbb{E}Y_{j_1}Y_{j_2}\dots Y_{j_{\pi}} = ( \mathbb{E} Y_{j_1})( \mathbb{E} Y_{j_2}) \dots (\mathbb{E}Y_{j_{|\pi|}}) = n^{-|\pi|}
\] 
so for fixed $\pi$
\[
    \sum_{J}^{*} \mathbb{E} Y_{j_1}\dots Y_{j_{|\pi|}} = o(1) \implies \sum_{\pi }^{ } \sum_{J}^{} \mathbb{E} \vec{Y}_J = O(1)
\] 
now we need
\[
\mathbb{E}(\sum_{i=1}^{n} X_i)^{k} - \mathbb{E}(\sum_{i=1}^{n}Y_i)^{k} = o(1)
\] 
LHS =
\[
\sum_{\pi }^{} \sum_{ J}^{*}\left( \mathbb{E}\vec X_J - \mathbb{E} \vec Y_J\right) = o(1)
\] 
if we can prove that the things in the sum are $o(n^{-|\pi|})$
\[
    \mathbb{E} Y_{j_1}\dots Y_{j_{|\pi|}} = n^{-|\pi|}
\] 
\[
\mathbb{E}\vec X_J = \mathbb{P}(X_{j_1} = X_{j_2}\dots X_{j_{|\pi|}} = 1) = \mathbb{P}(X_1 = X_2 = X_3 = \dots X_{|\pi|} = 1) 
\]
\[
= \frac{1}{n} \frac{1}{n-1} \frac{1}{n-2} \dots \frac{1}{n-|\pi| + 1} = \frac{1}{n^{|\pi|}} + O(\frac{1}{n^{|\pi| + 1}})
\] 

\section{Example weak CLT}
Let $\{X_k\}$ be independent random variables such that for all  $k, m$
\[
    \mathbb{E} X_{k} = 0 \;\; \mathbb{E} X_k^2 = 1 \;\;\; | \mathbb{E}X_k^{m}| \le C_m
\] 
want to show that
\[
    Z_n := \frac{\sum_{k=1}^{n}X_k}{\sqrt{n}} \stackrel{d}{\rightarrow} N(0,1)
\] 
First $N(0,1)$ has exponential decay, unique so need to prove  $ \mathbb{E}(Z_n)^{m} \rightarrow \mathbb{E}N(0,1)^{m}$ for all $m$.
Let $Y_k$ be i.i.d $N(0,1)$
 \[
\tilde Z_n : = \frac{\sum_{k=1}^{n}Y_k}{\sqrt{n}} \sim N(0,1)
\] 
we need to prove for all $m$
 \[
n^{-\frac{m}{2}}\left(\mathbb{E}(\sum_{k=1}^{n}X_k)^{m} - \mathbb{E}(\sum_{k=1}^{n}Y_k)^{m}\right) = o(1)
\] 
\[
    \mathbb{E}\left( \sum_{i=1}^{n}X_i\right)^{k} = \sum_{\pi }^{} \sum_{J}^{*} \mathbb{E} X_{i_1}\dots X_{i_k}|_{\pi,J}
\] 
but $X_1^2 \ne X_1$
but if $\pi = \{\pi_1,\pi_2,\pi_3,\dots,\pi_{|\pi|}\}$ then we get
\[
    \sum_{\pi }^{ } \sum_{J}^{*} \mathbb{E} X_{j_1}^{ |\pi_1|} X_{j_2}^{|\pi_2|} \dots X_{j_{|\pi|}}^{|\pi_\pi|} = \sum_{\pi }^{} \sum_{ J}^{*} \left( \mathbb{E} X_{j_1}^{|\pi_1|}\right) \dots \left( \mathbb{E}(X_{j_{|\pi|}}^{|\pi_{|\pi|}|} \right)
\] 
there are finitely many $\pi$ each $J $ is size  $n^{|\pi|}$ and $|\pi| \le k$ so the total is bounded by $\tilde C_k = O(1)$ since each moment is bounded so its the maximum moment times k. so for fixed $\pi$ we have $O(n^{|\pi|})$ $k= m$ since they were
used interchangably by the prof. now if we multiply by  $(n^{-\frac{k}{2}})$ we get for fixed $\pi$
\[
O(n^{|\pi| - \frac{k}{2}})
\] 
we only need to focus on the $\pi$'s such that $|\pi| \ge \frac{k}{2}$ \\
We use Lemma:\\
\emph{
    if $|\pi| < \frac{k}{2}$ there exists $\alpha$ such that $|\pi_\alpha| = 1$
}\\
so this implies that 
\[
    \mathbb{E}|X_{j_1}^{|\pi_1|} \dots  = 0
\] 
when $|\pi| > \frac{k}{2}$
 Now we only need to focus on the $\pi$'s such that
\[
|\pi| = \frac{k}{2}
\] 
but if we want to avoid singletons we only need to care about pairings, where each group is of size 2.


\section{Topic: Martingale}
Conditional expectation
\[
    \mathbb{E}[X | Y], \mathbb{E}[X| Y_1,Y_2,Y_3] 
\] 
\[
    \mathbb{E}[X | \mathcal{G}] \qquad \mathcal{G} \text{ is $\sigma$-algebra}
\] 
we define expectation first since we can recover conditional probability via
\[
    \mathbb{P}[A | \mathcal{G}] := \mathbb{E}[1_A | \mathcal{G}]
\] 
Review
\[
\mathbb{P}(A |B) = \frac{\mathbb{P}(A \cap B)}{ \mathbb{P}(B)}
\] 
where $ \mathbb{P}(B) \ne 0$.
where we extend $ \mathbb{E}[X | B] := \frac{\mathbb{E}[X 1_{B}]}{ \mathbb{P}(B)}$ 
and
\[
    \mathbb{E}[X | Y = y] = \mathbb{E}[X | \{\omega : Y(\omega) = y\}] = \frac{ \mathbb{E}X 1_{Y=y}}{ \mathbb{P}(Y =y)}
\] 
\[
    \mathbb{P}(X=x|Y=y) = \mathbb{P}_{X|Y}(x|y) = \frac{\mathbb{P}_{X,Y}(x,y)}{ \mathbb{P}_Y(y)}
\] 
\[
    \mathbb{E}[X | Y=y] = \sum_{k}^{} k \mathbb{P}_{X|Y}(k|y)
\] 
\[
    f_{X|Y}(x|y) = \frac{f_{X,Y}(x,y)}{f_Y(y)}
\] 
in 275+B we dont use
\[
    \mathbb{E}[X|Y=\frac{1}{2}]
\] 
when $Y=\frac{1}{2}$ has 0 probability
we use
\[
    \mathbb{E}[X|Y] := \mathbb{E}[X | \sigma(Y)]
\] 
\[
    \mathbb{E}X = \sum_{k}^{} \mathbb{E}[X | Y=k] \mathbb{P}(Y =k)
\] 
in math 170
\[
    \mathbb{E}[X|Y] = h(Y), h(k) := \mathbb{E}[X| Y=k]
\] 
we can write total expectation value formula
\[
    \mathbb{E}X =  \mathbb{E}[\mathbb{E}[X|Y]]
\] 
and conditional expectation value is a random variable.
\[
    \mathbb{E}[X|\mathcal{G}]
\] 

\section{Definition }
\emph{
    $\mathcal{G}$ is a  $\sigma$-algebra,  $\mathcal{G} \subset \mathcal{F}$ under $ (\Omega, \mathbb{P}, \mathcal{F})$ 
    and $ \mathbb{E}|X| < \infty$ then $ \mathbb{E}[X|\mathcal{G}]$ is defined as 
    a random variable $ \in \mathcal{G}$ i.e. $ \mathbb{E}[X|\mathcal{G}] \in \mathcal{G}$
}

recall $z \in \mathcal{G}, Z (\Omega, \mathcal{F}) \rightarrow ( \mathbb{R}, \mathcal{R})$ then $Z^{-1}(\mathcal{R}) \subset \mathcal{G}$.
For example if $Z \in \sigma(Y) \implies Z = h(Y)$ so $ \mathbb{E}[X | Y] = h(Y)$ for some $h$ measurable.
Example $\Omega = \mathbb{R}, \mathcal{F} =\mathcal{R}$ and $\mathcal{G} = \sigma(\bigcup_{n \in \mathbb{Z}} \{[n,n+1)\})$.
THen all possible elements in $\mathcal{G}$ are unions of unit intervals.
If $X \in  \mathcal{G}$ then  $X$ is constant on each interval


\section{Friday}

We want
\[
    \mathbb{E} X Y = \mathbb{E}( \mathbb{E}[X | \mathcal{G}] Y)
\] 
\[
= \mathbb{E}\tilde X Y
\] 

\emph{
    Def: $Y = 1_A, A \in \mathcal{G}$
}\\
Definition of $ \mathbb{E}[X | \mathcal{G}]$:\\
\emph{
    Let $\mathcal{G} \subset \mathcal{F}$, $\mathcal{G}$ is  $\sigma$-algebra $X$ is and  $ \mathbb{E}|X| < \infty$.
    If $\tilde X \in \mathcal{G}$ , we have for all  $A \in \mathcal{G}$
     \[
    \mathbb{E} X 1_A = \mathbb{E} \tilde X 1_A
    \] 
    then we say $\tilde X = \mathbb{E}[X | \mathcal{G}]$
}

First uniquenes, $B \in \mathcal{G}$ and  $ \mathbb{P}(B) = 0$
\[
\hat X = \tilde X + 2 1_B
\] 
and
\[
    \tilde X \in \mathcal{G} \implies \hat X \in \mathcal{G}
\] 
and
\[
\mathbb{E} \hat X 1_A = \mathbb{E} \tilde X 1_a + 2 \mathbb{E}1_b 1_a = \mathbb{E} X 1_a
\] 
proof of uniqueness.
Lemma: Let's say $Y_1,Y_2 \in \mathcal{G}$  and for any $A \in \mathcal{G}$  $ \mathbb{E}Y_1 1_a = \mathbb{E}Y_2 1_A = \mathbb{E}1_A X$ then
\[
\mathbb{P}(Y_1 \ne Y_2) = 0
\] 
Proof: Define $A_+ = \{Y_1 > Y_2\} \in \mathcal{G}$
\[
    \mathbb{E}1_{A^{+}}Y_1 = \mathbb{E}[1_{A^{+}}Y_2] \implies \mathbb{P}(A^{+}) = 0
\] 
this implies that
\[
\mathbb{P}(Y_1 \ne Y_2) = 0
\] 

\section{Example}
$\mathcal{G} = \mathcal{F}$ then $ \mathbb{E}[X | \mathcal{F}] = X] = X$

\end{document}
